\clearpage
\section{Logique}
\begin{objectif}
Repérer l'hypothèse d'une implication, construire sa négation et distinguer un témoin choisi pour chaque objet d'un témoin commun à tous.
\end{objectif}
\subsection{Explorer : un contrôle de fabrication}
Un capteur produit trois informations : $P$ signifie « pression conforme », $Q$ « température conforme », $R$ « contrôle manuel effectué ». Une pièce est acceptée si $(P\wedge Q)\vee R$. Dans le cube du diaporama, examiner les huit états possibles. Proposer un état où la pièce est acceptée malgré une pression non conforme. Décrire ensuite exactement les états refusés.
\subsection{Connecteurs et preuve}
\begin{definition}[valeur logique]
Une assertion mathématique possède une valeur vraie ou fausse. $\neg P$ inverse celle de $P$. $P\wedge Q$ exige les deux valeurs vraies ; $P\vee Q$ en exige au moins une. L'implication $P\Rightarrow Q$ est fausse précisément lorsque $P$ est vraie et $Q$ fausse. Une équivalence exige que les deux valeurs coïncident.
\end{definition}
\begin{propriete}[négation et contraposée]
On a $\neg(P\Rightarrow Q)\iff P\wedge\neg Q$ et $(P\Rightarrow Q)\iff(\neg Q\Rightarrow\neg P)$.
\end{propriete}
\begin{demonstration}
La seule situation qui contredit l'implication est $(P,Q)=(\text{vrai},\text{faux})$. C'est exactement la situation décrite par $P\wedge\neg Q$, et aussi la seule qui contredit $\neg Q\Rightarrow\neg P$.
\end{demonstration}
Une réciproque échange l'hypothèse et la conclusion. Elle nécessite une preuve séparée. Un contre-exemple doit satisfaire l'hypothèse et contredire la conclusion.
\subsection{Quantifier sans ambiguïté}
\begin{definition}[témoins]
$\forall x\in E,\ P(x)$ exige $P$ pour chaque élément de $E$. $\exists x\in E,\ P(x)$ demande au moins un témoin. La négation du premier énoncé est $\exists x\in E,\ \neg P(x)$ ; celle du second est $\forall x\in E,\ \neg P(x)$.
\end{definition}
Dans $\forall x\,\exists y$, le choix de $y$ peut dépendre de $x$. Dans $\exists y\,\forall x$, le même $y$ doit fonctionner partout. Pour réfuter cette dernière assertion, expliquer comment mettre en défaut chaque candidat $y$.
\subsection{Ateliers autonomes}
\begin{enumerate}
\item \textbf{Diagnostic.} Énumérer les états acceptés du contrôle de fabrication. Le responsable affirme : « toute pièce acceptée satisfait $P$ ». Décider si cette règle est correcte et écrire un état qui tranche la question.
\item \textbf{Un budget commun ?} Pour $x,y\in[2,6]$, comparer $\forall x\,\exists y,\ x+y=8$ et $\exists y\,\forall x,\ x+y=8$. Donner une preuve ou une réfutation, puis écrire les deux négations.
\item \textbf{Une preuve par les restes.} Un entier est divisible par 3 lorsqu'il s'écrit $3k$ avec $k\in\Z$. Montrer que si $m^2$ est divisible par 3, alors $m$ l'est aussi, en étudiant les trois restes possibles de $m$ dans une division par 3.
\end{enumerate}
\ifprof
\subsection{Préparer et discuter}
Faire distinguer une panne réelle d'une hypothèse fausse dans l'implication. Au deuxième atelier, deux étudiants jouent respectivement les rôles de celui qui choisit $x$ et de celui qui choisit $y$ ; inverser l'ordre des choix. Pour le dernier, nommer la contraposée avant de commencer le calcul.
\subsection{Solutions des ateliers}
\begin{enumerate}
\item En codant faux par 0 et vrai par 1, les états acceptés sont $(0, 0, 1), (0, 1, 1), (1, 0, 1), (1, 1, 0), (1, 1, 1)$. L'état $(0,0,1)$ réfute la règle : le contrôle manuel suffit à accepter, même lorsque P est fausse. Les états refusés satisfont $\neg R\wedge(\neg P\vee\neg Q)$.
\item La première assertion est vraie : choisir $y=8-x$, qui reste dans $[2,6]$. La seconde est fausse : $x=2$ imposerait $y=6$, tandis que $x=6$ imposerait $y=2$. Leurs négations sont respectivement $\exists x\in[2,6],\ \forall y\in[2,6],\ x+y\ne8$ et $\forall y\in[2,6],\ \exists x\in[2,6],\ x+y\ne8$.
\item Écrire $m=3q+r$ avec $r\in\{0,1,2\}$. Alors $m^2=9q^2+6qr+r^2$ a le même reste que $r^2$, soit 0, 1, 1 pour les trois cas. Si le reste du carré est nul, r doit être nul. C'est la contraposée : un entier non divisible par 3 a un carré non divisible par 3.
\end{enumerate}

\fi
